\section{Raw verification for inaccessible parameter families}\label{sec:uniform_realizations}
The following verifications use the unnormalized physical instruments. In particular, an unstable normalized posterior does not invalidate the comparison of all fixed programmes. The physical model is supplied for offline synthesis; its actual parameter and evolving state are not observer registers.

\begin{theorem}[Unknown, possibly nonmixing Gaussian models]\label{thm:uniform_gaussian}
Let $\Theta$ be computably compact. Consider the Gaussian experiment of Theorem~\ref{thm:effective_gaussian}, with uniformly computable preparations, transition matrices $T_{\theta,t}^a$, means $\mu_{\theta,t,j}^a$, and covariance matrices $\Sigma_{\theta,t,a}$. Suppose the same known bounds $U,\sigma_-,\sigma_+$ hold throughout the family. For parameter distance at most $s$, assume supplied uniform moduli
\[
 \|T_{\theta,t}^a-T_{\theta',t}^a\|_{1\to1}\le\tau_t(s),\quad
 \|\mu_{\theta,t,j}^a-\mu_{\theta',t,j}^a\|_2\le m_t(s),\quad
 \|\Sigma_{\theta,t,a}-\Sigma_{\theta',t,a}\|_F\le v_t(s).
\]
For categorical Brier prediction, or a uniformly presented bounded terminal task, Theorem~\ref{thm:uniform} applies, with the common report reduction \eqref{eq:gaussian_effective} and parameter-instrument modulus
\begin{equation}\label{eq:gaussian_parameter}
 \omega_t^I(s)=\tau_t(s)+\frac{m_t(s)}{\sigma_-}
                      +\frac{v_t(s)}{\sqrt2\,\sigma_-^2}.
\end{equation}
Thus the exact minimax risk over all parameter-independent finite-state observers has computably convergent upper and lower certificates at every fixed feasible $(n,M)$. Transitions may be reducible, deterministic, or permutations; no normalized-filter contraction, positive transition floor or full-rank predictive geometry is assumed here.
\end{theorem}
\begin{proof}
For a normalized column vector $x$, the instrument is
$\operatorname{diag}(\varphi_{\theta,t,j}^a(y))T_{\theta,t}^a x\,dy$.
Subtract first the transition matrix and then the emission density. Integration bounds the first term by $\tau_t(s)$ and the second by the largest $L^1$ difference of a pair of Gaussian densities. To bound that difference, interpolate the means and covariances on a line segment. Every intermediate covariance has eigenvalues at least $\sigma_-^2$. For a mean direction $b$, the $L^1$ norm of the density derivative is
\[
 \E|b^T\Sigma^{-1}(Y-\mu)|\le\|b\|_2/\sigma_-.
\]
For a covariance direction $C$, set $A=\Sigma^{-1/2}C\Sigma^{-1/2}$ and let $Z$ be standard Gaussian. The derivative norm is bounded by
\[
 \frac12\E|Z^TAZ-\tr A|
 \le\frac12\sqrt{2\tr(A^2)}
 \le\frac{\|C\|_F}{\sqrt2\,\sigma_-^2}.
\]
The middle identity follows by orthogonally diagonalizing $A$ and using independence and $\E(Z_i^2-1)^2=2$. Integrating the two directional bounds proves \eqref{eq:gaussian_parameter}.

All emission mixtures are everywhere positive for every parameter and normalized incoming state. The same box cells, one genuine tail cell, and state- and parameter-independent reconstruction measures therefore have the same support table. The uniform gradient and tail estimates in Theorem~\ref{thm:effective_gaussian} give \eqref{eq:instrument_error} without conditioning away the tail. Uniform computability and the displayed moduli give certified cell integrals at every point of a parameter net. Restriction to the future-test effect space gives the parameterwise predictive quotients, even when some states are predictively identical. Categorical Brier loss itself is independent of the inaccessible model, so $\omega^L=0$ in that case. Preparation variation is charged by $\omega_0$. The common-programme theorem now applies directly. All memory and timing constraints refer to the one synthesized tuple, not a different filter chosen at each parameter.
\end{proof}

\begin{theorem}[Unknown noncommuting channels and preparations]\label{thm:uniform_quantum}
On $\mathbb C^D$, with $D$ fixed, consider a computably compact family of instruments
\[
 \mathsf I_{\theta,t}^a(dy)\rho
   =E_{\theta,y}^{1/2}\mathcal C_{\theta,t}^a(\rho)
                                  E_{\theta,y}^{1/2}\,\zeta(dy),
 \qquad \int E_{\theta,y}\,\zeta(dy)=I,
\]
where $\mathcal C_{\theta,t}^a$ is completely positive and trace preserving,
$\kappa I\preceq E_{\theta,y}\preceq K I$, and $\zeta$ is a fixed compact-report reference probability. Assume uniform effective report partitions and instrument oscillation bounds, as in Theorem~\ref{thm:effective_quantum} or Proposition~\ref{prop:effective_singular}. For nearby parameters suppose
\[
 \sup_{\rho}\|\mathcal C_{\theta,t}^a(\rho)-\mathcal C_{\theta',t}^a(\rho)\|_1
          \le c_t(s),\qquad
 \sup_y\|E_{\theta,y}-E_{\theta',y}\|_{\rm op}\le e_t(s),
\]
with supplied computable moduli. Then Theorem~\ref{thm:uniform} applies to every uniformly presented bounded terminal task, in particular Brier prediction of a fixed finite POVM, with
\begin{equation}\label{eq:quantum_parameter}
              \omega_t^I(s)=c_t(s)+\sqrt{K/\kappa}\,e_t(s).
\end{equation}
Initial preparations may be pure and channels may be unitary. The effects need not commute. The persistent resource is the classical observer alphabet; the experimental quantum state is not free quantum memory of that observer.
\end{theorem}
\begin{proof}
Replace the channel output first. For a Hermitian difference $v$, positivity of the instrument and normalization of its integrated trace imply
\[
 \int\|E_{\theta,y}^{1/2}vE_{\theta,y}^{1/2}\|_1\,\zeta(dy)\le\|v\|_1.
\]
Apply the instrument separately to the positive and negative parts of $v$ to obtain this inequality, so the channel contribution is at most $c_t(s)$, not a division by a posterior eigenvalue. For the second replacement, the Sylvester-equation estimate proved in Theorem~\ref{thm:effective_quantum} gives
$\|E_{\theta,y}^{1/2}-E_{\theta',y}^{1/2}\|_{\rm op}
 \le e_t(s)/(2\sqrt\kappa)$.
Expand the two-factor congruence difference at a density matrix. Its trace norm is at most $\sqrt{K/\kappa}\,e_t(s)$; integrating proves \eqref{eq:quantum_parameter}.

Every report density is between $\kappa$ and $K$ on the same reference support, irrespective of the parameter or purity of the state. Common reconstruction therefore preserves exact supports. Computable positive square roots, channel entries, cell integrals and decision nets give the effective data. Finite-dimensional closure of future-test effects supplies each predictive quotient; indistinguishable decompositions of a density matrix are not retained as different predictive points. Apply Theorem~\ref{thm:uniform}. For a unitary family, for example, $c_t(s)\le2\sup\|U_{\theta,t}^a-U_{\theta',t}^a\|_{\rm op}$, by expanding the conjugation difference. No reset or mixing is needed.
\end{proof}

\begin{corollary}[Singular acquisition with parameter uncertainty]\label{cor:uniform_singular}
Suppose a parameter family has a common Cantor reference, uniformly positive instrument densities with a supplied uniform H\"older modulus, and parameter variation bounded as in \eqref{eq:parameter_moduli}. Ternary cylinders give the same $d_{t,k}\le L_t3^{-\alpha k}$ for all parameters. Likewise a common countable support with a computable uniform actual tail bound $T_J\downarrow0$ gives $d_{t,J}\le2T_J$. Separately flagged finite instrument branches may have different output ranks, provided their active probabilities and report supports satisfy the common support condition. In each case \eqref{eq:uniform_interval} holds at unchanged $M$.
\end{corollary}
\begin{proof}
Use the same conditional reference on each cylinder or the same tail reconstruction at every parameter. The proof of Proposition~\ref{prop:effective_singular} bounds the positive vector-measure error by the cylinder oscillation or twice the actual tail mass, uniformly. The new parameter modulus is an additional bound on the unnormalized instruments, not a claim about Hausdorff dimension. Flagged branches retained exactly add no binning error. Apply the common-programme reduction and Theorem~\ref{thm:uniform}.

For an explicit rank-changing family, take the Cantor qubit instrument in Proposition~\ref{prop:effective_singular}, allow its coefficient to range over a known compact subinterval of $(0,1)$, and mix it with separately reported pure- and full-rank preparation branches whose weights range over a compact positive simplex. The branch maps are positive, their weights are bounded away from zero, and their coefficient/weight dependence supplies \eqref{eq:parameter_moduli} directly. At the first call, the acquired law is a genuine mixture of an injective affine image of the Cantor law and atoms of different ranks. The original branch masses enter every law and every bound. Subsequent masses are the actual history-dependent ones; no singular stratum is normalized to probability one. This is a raw family with both inaccessible coefficients and changing acquired rank, not merely a declaration of a singular interface.
\end{proof}

\subsection{An exact calibration--memory threshold}
The following small experiment makes the common-parameter quantifier and the price of retaining acquired calibration information explicit. It is a sharp consequence, not a claim to resolve a general nonlinear filtering problem.

Choose an unknown fixed bit $\vartheta\in\{0,1\}$ and an unknown fixed noise level $p\in[p_-,p_+]$, where $0<p_-\le p_+<1/2$. The first paid call reports
$C=\vartheta\oplus E$, with $E\sim\operatorname{Bernoulli}(p)$. The second paid call reports an independent fair bit $Y$. At halt the observer forecasts $X=\vartheta\oplus Y$ under scalar squared loss. A post-decision audit scores that forecast and reveals no information before it is made. Thus there are two observer acquisitions and, if scoring is charged as an acquisition, one additional audit call. There is one action at both phases and the same binary report alphabet without a phase flag. One programme must serve the whole family. This is a finite positive raw kernel with nonconstant acquired calibration information, and all report supports coincide.

\begin{proposition}[A shared five-state implementation and its sharp converse]\label{prop:uniform_threshold}
For the preceding task, external widths $(m_1,m_2)$ satisfy
\begin{equation}\label{eq:threshold_external}
 R^\Theta_{\ext}(m_1,m_2)=
 \begin{cases}
 p_+(1-p_+),&m_1,m_2\ge2,\\
 1/4,&\min(m_1,m_2)=1.
 \end{cases}
\end{equation}
The internally exact two-call class is empty for $M<3$ and has
\begin{equation}\label{eq:threshold_aut}
 R^\Theta_{\aut}(2,M)=
 \begin{cases}1/4,&3\le M\le4,\\ p_+(1-p_+),&M\ge5.\end{cases}
\end{equation}
These assertions allow all randomized acquisition and transition rows of the respective classes. If an additional unobserved fixed offset $b\in[-\delta,\delta]$ is added to the audited target, both nonempty risk formulas increase by exactly $\delta^2$. With the true bit furnished to the programme in advance, the unshifted risk is zero at $M=4$; furnishing it is a strictly different information pattern.
\end{proposition}
\begin{proof}
For a lower bound even on full-history observers, fix $p=p_+$ and put a uniform prior on $\vartheta$. Given $(C,Y)$ the target equals $C\oplus Y$ with conditional probability $1-p_+$ and its complement with probability $p_+$. Conditional squared-loss projection gives minimum Bayes risk $p_+(1-p_+)$. A worst-case risk is at least this average. Conversely, retain $C$ and then retain $W=C\oplus Y$. Output $p_+$ at $W=0$ and $1-p_+$ at $W=1$. For every $\vartheta$ and $p$ its risk is
\[
 (1-p)p_+^2+p(1-p_+)^2
      =p_+^2+p(1-2p_+)\le p_+(1-p_+).
\]
This is one common programme; it proves the upper with two labels at both cuts.

If $m_1=1$, the second-call forecast is downstream of $Y$ and fresh randomness, independent of the erased $C$. Under the uniform prior on $\vartheta$, the target is then a fair bit conditional on the available information, giving risk at least $1/4$. If $m_2=1$, the fixed terminal forecast sees neither bit, and $X$ is fair even for each fixed $\vartheta$. The constant $1/2$ attains the bound. Terminal randomization cannot improve it, by convexity of squared loss. These arguments prove \eqref{eq:threshold_external}.

At either phase both reports have positive probability under each model. The common-null-set clock argument of Proposition~\ref{prop:clock} forces the initial, intermediate and stopping label supports to be disjoint. They are also common across parameters, by positivity and the common programme. Their cardinalities therefore obey $1+m_1+m_2\le M$, giving emptiness below three and a one-label cut when $M\le4$. Five states suffice: an initial label $s$, intermediate labels $c_0,c_1$, and stopping labels $d_0,d_1$. For the single legal action, the common update and readout are
\[
 k(c_i\mid s,y)=\boldsymbol1_{\{i=y\}},\qquad
 k(d_j\mid c_i,y)=\boldsymbol1_{\{j=i\oplus y\}},\qquad
 u(d_0)=p_+,\quad u(d_1)=1-p_+.
\]
All other occupied-row entries vanish. The labels themselves implement the deadline. The constant-risk three-state programme is a chain of one initial, one intermediate and one stopping label with readout $1/2$. These exact readout values belong to the declared mathematical programme; finite-bit evaluation of a nonrational $p_+$ is charged by its output precision, not by an assumption that the true $p$ is known. This proves \eqref{eq:threshold_aut}, including the converse for all randomized rows.

For the offset extension, place the independent uniform prior on $b\in\{-\delta,\delta\}$ in the lower argument. It adds conditional variance $\delta^2$, while the one-label lower arguments remain valid. In the upper constructions, the mean forecast error is zero for each $\vartheta,p$, since $Y$ is fair. Thus the cross term with every fixed $b$ vanishes and the risk increases by $b^2$, at most $\delta^2$. For an oracle knowing $\vartheta$, ignore $C$ at a single intermediate state and encode $\vartheta\oplus Y$ in two stopping states, using four states in total. Its zero risk does not give a feasible robust programme. This last comparison explains why the parameter supremum must follow the common-programme choice in \eqref{eq:robust_risk}.
\end{proof}


\begin{corollary}[Attained calibration and simulation coordinates]\label{cor:uniform_erasure}
In Proposition~\ref{prop:uniform_threshold}, erase the first report independently with a fixed known probability $a\in(0,1)$, without making an extra observation. The second report remains the independent fair bit. Give the simulator one physical first-report call and require it to complete the virtual first report before the second report or the postdecision audit. The exact parameter-uniform causal deficiency to the unerased experiment is
\begin{equation}\label{eq:uniform_exact_deficiency}
 \varepsilon=a(1/2-p_-).
\end{equation}
For every $M\ge7$, the worst-case squared prediction risk in this erased experiment, with the same unobserved offset $b\in[-\delta,\delta]$, is exactly
\begin{equation}\label{eq:uniform_all_coordinates}
 R^{\Theta}_{\aut}(2,M,\delta)
 =p_+(1-p_+)
  +\frac{(1/2-p_+)^2}{1/2-p_-}\,\varepsilon+\delta^2.
\end{equation}
The optimum uses one parameter-independent seven-state programme. There are two observer acquisitions; an optional scored physical audit is a third call after the decision and does not enlarge the observer deadline. No assertion about the optimal number of labels below seven is made for this enlarged report alphabet.
\end{corollary}
\begin{proof}
On a nonerased first report the simulator copies the bit. On an erasure it supplies a fresh fair bit. It copies the second bit. This is one report-to-report kernel, with no retained information, for all parameters and phases. At a fixed $(\vartheta,p)$, the virtual error probability is $(1-a)p+a/2$, so its total variation distance from the target first-report law is $a(1/2-p)$. The second report is independent, and the target scored variable is determined by that report and the fixed parameter. Thus the same bound holds for the complete stopped and scored law. Taking the supremum gives the upper bound in \eqref{eq:uniform_exact_deficiency}.

For the converse fix $p=p_-$ and compare $\vartheta=0$ and $\vartheta=1$. The target first-report laws have total variation distance $1-2p_-$; the source laws have distance $(1-a)(1-2p_-)$. Every common first-report simulator contracts total variation. If its error in each model is at most $e$, the triangle inequality therefore gives
\[
 1-2p_-\le (1-a)(1-2p_-)+2e.
\]
Consequently $e\ge a(1/2-p_-)$. No second report or audit is available when this virtual report is due. This proves the claimed exact deficiency and explains the importance of the interface.

For a lower risk bound put the uniform prior on $\vartheta$, fix $p=p_+$, and put an independent uniform prior on $b\in\{-\delta,\delta\}$. Even with the complete two-report history, the conditional variance of $X$ is $p_+(1-p_+)$ when the first bit survives and $1/4$ when it is erased. The offset contributes $\delta^2$. Hence every programme has worst-case risk at least
\[
 (1-a)p_+(1-p_+)+a/4+\delta^2
 =p_+(1-p_+)+a(1/2-p_+)^2+\delta^2.
\]
For attainment, use one initial state, three intermediate states recording $0,1,*$, and three terminal states with readouts $p_+,1-p_+,1/2$. On a surviving bit, update by the same XOR rule as in Proposition~\ref{prop:uniform_threshold}; on an erasure, send either second report to the last terminal state. This is a fixed two-step programme, with all seven states charged. At each fixed $(\vartheta,p)$ its risk without the offset is
\[
 (1-a)\{p_+^2+p(1-2p_+)\}+a/4.
\]
It increases with $p$ and has zero mean forecast error, since the second bit is fair. Every fixed offset therefore adds exactly $b^2$. The worst-case risk equals the lower bound. Substitute \eqref{eq:uniform_exact_deficiency} to obtain \eqref{eq:uniform_all_coordinates}.

The erasure flag has phase-dependent support, although its support is common across inaccessible parameters at each phase. Thus this construction meets the support-resolved uniform theorem, but the common-null-set clock converse for the binary-report experiment cannot simply be reused to claim a new sharp threshold below seven. The matching risk in \eqref{eq:uniform_all_coordinates} comes from a full-history lower bound and an explicit feasible programme, not such a converse.
\end{proof}
